\begin{pastexam}[label={px:ncku-114-prog-01}]{成大資訊聯招 114 學年度 第 1 題}
（單選）

{[3\%]} Consider the following two functions written in C language:
\begin{minted}[fontsize=\small]{c}
void enqueue(int queue[], int *front, int *rear, int size, int value) {
    if ((*rear + 1) % size == *front) {
        printf("Queue is full\n");
        return;
    }
    if (*front == -1) *front = 0;
    *rear = (*rear + 1) % size;
    queue[*rear] = value;
}
int dequeue(int queue[], int *front, int *rear, int size) {
    if (*front == -1) {
        printf("Queue is empty\n");
        return -1;
    }
    int value = queue[*front];
    if (*front == *rear) {
        *front = *rear = -1;
    } else {
        *front = (*front + 1) % size;
    }
    return value;
}
\end{minted}
Assume the queue Q with size = 5, initially empty (front = -1, rear = -1). What will be the content of the array Q after executing the given sequence of operations?
\begin{minted}[fontsize=\small]{c}
enqueue(Q, &front, &rear, size, 10);
enqueue(Q, &front, &rear, size, 20);
dequeue(Q, &front, &rear, size);
enqueue(Q, &front, &rear, size, 30);
enqueue(Q, &front, &rear, size, 40);
enqueue(Q, &front, &rear, size, 50);
dequeue(Q, &front, &rear, size);
enqueue(Q, &front, &rear, size, 60);
enqueue(Q, &front, &rear, size, 80);
enqueue(Q, &front, &rear, size, 90);
dequeue(Q, &front, &rear, size);
dequeue(Q, &front, &rear, size);
enqueue(Q, &front, &rear, size, 100);
\end{minted}
\begin{enumerate}[label=(\Alph*)]
\item \{40, 50, 60, 80, 90\}
\item \{80, 90, 100, 50, 60\}
\item \{-1, 50, 60, 80, 100\}
\item \{50, 60, 80, 100, 40\}
\item \{60, 80, 100, -1, 50\}
\end{enumerate}
\end{pastexam}

\begin{pastsoln}
答案：(D)。

\textbf{程式在做什麼}（資結書第~6.2.2~節的環狀佇列）。 \texttt{front} 指向第一個元素，\texttt{rear} 指向最後一個
元素，佇列空時兩者都是 $-1$。 \texttt{enqueue} 在 \texttt{(rear + 1) \% size == front} 時判定為滿，所以大小 $5$
的陣列最多放 $5$ 個元素；\texttt{dequeue} 只移動 \texttt{front}，\textbf{不會清掉}陣列裡的值。 題目問「陣列
\texttt{Q} 的內容」，所以已經出佇列的值仍然留在陣列裡。

\textbf{逐步追蹤}（陣列寫 \texttt{Q[0..4]}，空格寫 $\cdot$）：
\begin{center}
\small
\begin{tabular}{lccl}
\toprule
操作 & \texttt{front} & \texttt{rear} & \texttt{Q[0..4]} \\
\midrule
enqueue 10 & 0 & 0 & $10, \cdot, \cdot, \cdot, \cdot$ \\
enqueue 20 & 0 & 1 & $10, 20, \cdot, \cdot, \cdot$ \\
dequeue（取出 10） & 1 & 1 & $10, 20, \cdot, \cdot, \cdot$ \\
enqueue 30 & 1 & 2 & $10, 20, 30, \cdot, \cdot$ \\
enqueue 40 & 1 & 3 & $10, 20, 30, 40, \cdot$ \\
enqueue 50 & 1 & 4 & $10, 20, 30, 40, 50$ \\
dequeue（取出 20） & 2 & 4 & $10, 20, 30, 40, 50$ \\
enqueue 60 & 2 & 0 & $60, 20, 30, 40, 50$ \\
enqueue 80 & 2 & 1 & $60, 80, 30, 40, 50$ \\
enqueue 90 & 2 & 1 & 滿（$(1 + 1) \bmod 5 = 2 = \texttt{front}$），印出 Queue is full \\
dequeue（取出 30） & 3 & 1 & $60, 80, 30, 40, 50$ \\
dequeue（取出 40） & 4 & 1 & $60, 80, 30, 40, 50$ \\
enqueue 100 & 4 & 2 & $60, 80, 100, 40, 50$ \\
\bottomrule
\end{tabular}
\end{center}
最後陣列是 \texttt{Q[0..4]} $= 60, 80, 100, 40, 50$，\texttt{front} $= 4$、\texttt{rear} $= 2$，佇列裡由前到後是
$50, 60, 80, 100$；\texttt{Q[3]} 的 $40$ 已經出佇列，只是沒被清掉。

\textbf{對選項}。 五個選項都不是照索引 $0$--$4$ 寫的 $60, 80, 100, 40, 50$。 從 \texttt{front} 開始繞一圈讀陣列：
\texttt{Q[4], Q[0], Q[1], Q[2], Q[3]} $= 50, 60, 80, 100, 40$，正是 (D)，前四個是佇列裡的元素，最後的 $40$ 是殘留的
舊值。 其他選項都不可能：
\begin{itemize}
\item (A) 含 $90$：$90$ 入佇列時佇列已滿，沒有放進去。
\item (B) 也含 $90$，而且沒有 $40$。
\item (C)、(E) 含 $-1$：$-1$ 只是 \texttt{front}、\texttt{rear} 表示「空」的值，從來沒有寫進陣列。
\end{itemize}
所以選 (D)。 題目的說法不夠精確（陣列本身是 $60, 80, 100, 40, 50$），作答時可以把兩種讀法都寫出來。
\end{pastsoln}
